\subsection{GROUNDING: NLI-based support filtering}\label{subsec:knowledge-grounding}

The GROUNDING stage filters each MAP finding, by checking if the cited text actually entails the claim. This text is not the possibly paraphrased version generated by the LLM as evidence, but the text that is found after a database lookup, using \texttt{te\_id} in the output. PubMedBERT, a biomedical NLI cross-encoder, fine-tuned on MNLI and MedNLI \parencite{gu2021biomedbert, williams2018mnli, romanov2018mednli, reimers2019sbert}, is used to determine whether the source text entails the claim: the model outputs an entailment score between 0--1, and findings whose entailment score is below a certain threshold are dropped. 

GROUNDING is a supporting filter, not a correctness check: it removes findings with claims that appear not to be supported by their evidence, but a surviving claim is not guaranteed to be scientifically correct -- the output is never verified by a domain expert. 

The same NLI model is reused in the RELATE stage (Subsection~\ref{subsec:knowledge-relate}). In GROUNDING, it scores a source paragraph against an extracted claim; in RELATE, it scores one rule against another. The two stages use the shared NLI backbone differently: GROUNDING applies one entailment threshold to decide whether a finding should be kept or discarded, whereas RELATE applies separate entailment and contradiction thresholds to classify each canonical rule pair as supporting, contradicting or unrelated. The empirical behavior of the shared backbone is reported in Chapter~\ref{chapter:Results}, and the results are discussed in Chapter~\ref{chapter:Discussion}.